\documentclass[10pt,a4paper]{amsart}
\usepackage[margin=19mm]{geometry}
\usepackage{amsmath,amssymb,mathtools}
\usepackage[scr=rsfs,cal=euler]{mathalfa}
\usepackage{xfrac}
\usepackage[unicode,hidelinks,bookmarks=false]{hyperref}
\newcommand{\ie}{i.e.\ }
\newcommand{\cf}{cf.\ }
\newcommand{\resp}{respectively\ }
\newcommand{\Subsection}{\subsection}
\providecommand{\nobreakdash}{\nobreak\hbox{-}}
\setlength{\emergencystretch}{2em}
\multlinegap0pt
\usepackage{amssymb}
\usepackage{xcolor}
\newtheorem{theorem}{Theorem}
\newtheorem{proposition}{Proposition}
\newtheorem{corollary}{Corollary}
\newcommand{\navierbc}{\left[(2S({u}))\nn + \alpha {u}\right]_{\tanproj} = 0}
\newcommand{\nn}{{n}}
\newcommand{\tanproj}{\text{tan}}
\title{Exponential decay for the 3D Boussinesq equations with Navier boundary conditions}
\author{Wen Feng, Weinan Wang, Xiang Xu}
\date{Source: arXiv:2608.00972v1 (CC BY 4.0)}
\begin{document}
\maketitle

\noindent\emph{Source and licence.} Exponential decay for the 3D Boussinesq equations with Navier boundary conditions. Version arXiv:2608.00972v1 (CC BY 4.0), distributed by arXiv under the Creative Commons Attribution 4.0 International licence, http://creativecommons.org/licenses/by/4.0/, as recorded in the arXiv licence metadata for this exact version and shown on the abstract page https://arxiv.org/abs/2608.00972v1. Full licence notice: \texttt{licenses/arxiv-2608.00972-CC-BY-4.0.txt}.\par\smallskip

\noindent\textit{Editorial scope.} Counted unit: the article's theorem thm:main (source0.tex lines 1074--1079). The article's complete original proof of it is delivered with it (source0.tex lines 1081--1136, one \texttt{proof} environment of 56 lines). No mathematics is written, repaired or re-derived: the statement and the proof are the article's own lines, in the article's own notation, and no retained line is substituted or re-flowed.  Retained uncounted as the source-local context the proof needs: lines 194--198; lines 204--206; lines 849--871; lines 1000--1002; lines 1353--1372; lines 1426--1439.  Nothing was omitted from any retained span, so the package carries no omission record.  No retained line contains a citation, so the wrapper carries no bibliography and no bibliography entry.  No retained line contains a figure environment, an image-inclusion command or drawing code, so no image is delivered and the package declares no asset dependency.  Editorial formatting only, in addition: an AMS \texttt{amsart} preamble that restates the article's own declarations and macros used by the retained lines, namely the environments \texttt{theorem}, \texttt{proposition}, \texttt{corollary} on separate counters, together with the standard packages those lines need.  The retained environments print in this document as Theorem 1, Proposition 1, Corollary 1 in order of appearance, whereas the article prints the same items with the numbering of its own full document.  The complete native source of the article as distributed in the arXiv e-print for this exact version is bundled unchanged as \texttt{sources/206583/source0.tex}.
\begin{align}
    \partial_t {u} - \nu \Delta {u} + ({u} \cdot \nabla) {u} + \nabla p &= \rho {e}_3 \quad \text{in } (0, T) \times \Omega, \label{eq:bouss_momentum} \\
    \partial_t \rho - \kappa \Delta \rho + {u} \cdot \nabla \rho &= 0 \quad \text{in } (0, T) \times \Omega, \label{eq:bouss_density} \\
    \nabla \cdot {u} &= 0 \quad \text{in } (0, T) \times \Omega, \label{eq:bouss_div}
\end{align}

\begin{align}
    {u} \cdot \mathbf n = 0, \quad \navierbc, \quad \text{and} \quad \rho = 0 \quad \text{on } (0, T) \times \partial\Omega. \label{eq:navier_bc}
\end{align}

\begin{align}
&\frac{1}{2}\|u(t)\|_{L^2}^2
  + \nu \int_s^t \|\nabla u\|_{L^2}^2\,d\tau
\nonumber \\
&\qquad \leq
  \frac{1}{2}\|u(s)\|_{L^2}^2
  +\nu\int_s^t \int_{\partial\Omega}
  \bigl(\Lambda_\Omega-\alpha\bigr)
  |u_{\tanproj}|^2\,dS\,d\tau
  +\int_s^t\int_\Omega \rho e_3\cdot u\,dx\,d\tau,\label{eq:energy_ineq1}
\\[0.5em]
&\frac{1}{2}\|u(t)\|_{L^2}^2
  +2\nu\int_s^t\|S(u)\|_{L^2}^2\,d\tau
\nonumber \\
&\qquad \leq
  \frac{1}{2}\|u(s)\|_{L^2}^2
  -\nu\int_s^t\int_{\partial\Omega}
  \alpha|u_{\tanproj}|^2 \, dS \, d\tau
  +\int_s^t\int_\Omega \rho e_3\cdot u\,dx\,d\tau,\label{eq:energy_ineq2}
\\[0.5em]
&\frac{1}{2}\|\rho(t)\|_{L^2}^2
  +\kappa\int_s^t\|\nabla\rho\|_{L^2}^2\,d\tau\leq\frac{1}{2}\|\rho(s)\|_{L^2}^2.\label{eq:rho_energy_ineq}
\end{align}

\begin{equation}\label{kornpoincare2}
\|u\|^2_{L^2(\Omega)}\le C_{(KP,\alpha)}\left(\|S(u)\|_{L^2(\Omega)}^2+\int_{\partial\Omega}\alpha|u|^2\,dS\right)
\end{equation}

\begin{proposition}\label{AppendixA}
Let {$y$ be a fixed measurable representative of an
element of \(L^1_{\mathrm{loc}}([0,\infty))\)}, let
$q:[0,\infty)\to[0,\infty)$, and let $K,D',C'>0$ with $K\neq D'$.
{All pointwise inequalities below are understood for this
fixed representative.} Assume that there exists a full-measure set $\mathcal T\subset[0,\infty)$ with $0\in\mathcal T$ such that, for every $s\in\mathcal T$ and every $t\ge s$,
\begin{equation}\label{Aq}
q(t)\le q(s)e^{-D'(t-s)}
\end{equation}
and
\begin{equation}\label{A1}
y(t)+K\int_s^t y(\tau)\,d\tau
\le y(s)+C'q(s)\int_s^t e^{-D'(\tau-s)}\,d\tau.
\end{equation}
Then, for every $s\in\mathcal T$ and every $t\ge s$,
\begin{equation}\label{A2}
y(t)\le y(s)e^{-K(t-s)}+\frac{C'q(s)}{K-D'}\left(e^{-D'(t-s)}-e^{-K(t-s)}\right).
\end{equation}
No sign assumption on $y$ is required.
\end{proposition}

\begin{corollary}\label{Acoro}
Let {$y$ be a fixed measurable representative of an element
of \(L^1_{\mathrm{loc}}([0,\infty))\)} and let $K>0$.
{All pointwise inequalities below are understood for this
fixed representative.} Assume that there exists a full-measure set $\mathcal T\subset[0,\infty)$ with $0\in\mathcal T$ such that, for every $s\in\mathcal T$ and every $t\ge s$,
\[
y(t)+K\int_s^t y(\tau)\,d\tau\le y(s).
\]
Then
\[
y(t)\le y(s)e^{-K(t-s)}
\]
for every $s\in\mathcal T$ and every $t\ge s$.
\end{corollary}

\begin{theorem}[Exponential decay] \label{thm:main}
Let $(u,\rho)$ be a global Leray--Hopf weak solution to the Boussinesq system \eqref{eq:bouss_momentum}--\eqref{eq:navier_bc} with initial data $(u_0,\rho_0)\in\mathcal H$. Set $A_{\partial\Omega}:=\{x\in\partial\Omega:\alpha(x)>0\}$ and assume that $\mathcal H^2(A_{\partial\Omega})>0$. Then there exist constants $K,D>0$, depending only on $\nu$, $\kappa$, $\Omega$, and $\alpha$, such that 
\[
\| {u}(t)\|_{L^2}^2+\|\rho(t)\|_{L^2}^2\leq D\left(\|{u}_0\|_{L^2}^2+\|\rho_0 \|_{L^2}^2\right) e^{-Kt}.
\]
\end{theorem}

\begin{proof}
Let $\mathcal T$ be the full-measure set in the Leray--Hopf energy inequalities. The Poincar\'e inequality gives a constant $C_P=C_P(\Omega)>0$ such that, for almost every time,
\begin{equation}
    \begin{split}
        \|\rho\|^2_{L^2(\Omega)}\le C_P\|\nabla\rho\|^2_{L^2(\Omega)}.
    \end{split}
\end{equation} 
Applying it to \eqref{eq:rho_energy_ineq}, we obtain 
\begin{equation}
   \begin{split}
      \|\rho(t)\|^2_{L^2}+\frac{2\kappa}{C_P}\int_s^t\|\rho(\tau)\|^2_{L^2}d\tau\le \|\rho(s)\|^2_{L^2}
   \end{split}
\end{equation}
for every $s\in\mathcal T$ and every $t\ge s$. 
Applying Corollary~\ref{Acoro} on this same set $\mathcal T$, we obtain
\begin{equation}\label{rho1}
    \|\rho(t)\|^2_{L^2(\Omega)}\le\|\rho(s)\|^2_{L^2(\Omega)}e^{-\frac{2\kappa}{C_P}(t-s)}.
\end{equation}
In particular, taking $s=0$ gives
\begin{equation}\label{rho2}
    \|\rho(t)\|^2_{L^2(\Omega)}\le\|\rho(0)\|^2_{L^2(\Omega)}e^{-\frac{2\kappa}{C_P}t}.
\end{equation}

We next estimate the velocity. Applying \eqref{eq:energy_ineq2} for every $s\in\mathcal T$ and every $t\ge s$, \eqref{kornpoincare2}, Young's inequality, and \eqref{rho1}, we obtain 
\begin{equation}\label{energyu2}
\begin{split}
\|u(t)\|^2_{L^2(\Omega)}+&(\frac{2\nu}{C_{(KP,\alpha)}}-\epsilon)\int_s^t\|u(\tau)\|^2_{L^2(\Omega)}\,d\tau 
\\&\le \frac{1}{\epsilon}\int_s^t\|\rho(\tau)\|^2_{L^2{(\Omega)}}\,d\tau+\|u(s)\|^2_{L^2(\Omega)}\\& \le\frac{1}{\epsilon}\|\rho(s)\|^2_{L^2(\Omega)}\int_s^te^{-\frac{2\kappa}{C_P}(\tau-s)}\,d\tau+\|u(s)\|^2_{L^2(\Omega)}.
\end{split}
\end{equation}
Choose $\epsilon>0$ such that $\frac{2\nu}{C_{(KP,\alpha)}}-\epsilon>0$ and $\frac{2\nu}{C_{(KP,\alpha)}}-\epsilon\neq \frac{2\kappa}{C_P}$. Equations~\eqref{rho1} and \eqref{energyu2} verify, on the common set $\mathcal T$, the two hypotheses of Proposition~\ref{AppendixA}.  Applying that proposition yields
\begin{equation}\label{u0}
\begin{split}
    &\|u(t)\|^2_{L^2(\Omega)}\le e^{-(\frac{2\nu}{C_{(KP,\alpha)}}-\epsilon)(t-s)}\|u(s)\|^2_{L^2(\Omega)}\\
    +&\frac{1}{\epsilon}\|\rho(s)\|^2_{L^2(\Omega)}\big[\frac{e^{-\frac{2\kappa}{C_P}(t-s)}-e^{-(\frac{2\nu}{C_{(KP,\alpha)}}-\epsilon)(t-s)}}{\frac{2\nu}{C_{(KP,\alpha)}}-\epsilon-\frac{2\kappa}{C_P}}\big],
\end{split}
\end{equation}
and
\begin{equation}\label{u1}
    \|u(t)\|^2_{L^2(\Omega)}\le e^{-(\frac{2\nu}{C_{(KP,\alpha)}}-\epsilon)t}\|u(0)\|^2_{L^2(\Omega)}+\frac{1}{\epsilon}\|\rho(0)\|^2_{L^2(\Omega)}\big[\frac{e^{-\frac{2\kappa}{C_P}t}-e^{-(\frac{2\nu}{C_{(KP,\alpha)}}-\epsilon)t}}{\frac{2\nu}{C_{(KP,\alpha)}}-\epsilon-\frac{2\kappa}{C_P}}\big]
\end{equation}
when $s=0$.
Set
\[
K:=\min\left\{\frac{2\nu}{C_{(KP,\alpha)}}-\epsilon,\frac{2\kappa}{C_P}\right\},
\qquad
D:=1+\frac{1}{\epsilon\left|\frac{2\nu}{C_{(KP,\alpha)}}-\epsilon-\frac{2\kappa}{C_P}\right|}.
\]
Combining \eqref{rho2} and \eqref{u1}, we obtain
\[
\|u(t)\|_{L^2(\Omega)}^2+\|\rho(t)\|_{L^2(\Omega)}^2
\le
D\bigl(\|u_0\|_{L^2(\Omega)}^2+\|\rho_0\|_{L^2(\Omega)}^2\bigr)e^{-Kt}.
\]
This proves the result.
\end{proof}

\end{document}
